\documentclass[aspectratio=169]{beamer}
\usetheme{Madrid}
\usecolortheme{default}
\usepackage[utf8]{inputenc}
\usepackage[T5]{fontenc}
\usepackage{graphicx}
\usepackage{hyperref}
\usepackage{booktabs}
\usepackage{amsmath}

\setbeamertemplate{caption}[numbered]
\renewcommand{\figurename}{Hình}
\renewcommand{\thefigure}{19.\arabic{figure}}

\title[Chương 19: The Future of AI]{Học Sâu (Deep Learning) \\ \vspace{0.5cm} \Large Chương 19: Tương lai của Trí tuệ Nhân tạo (The Future of AI)}
\author{Đại học Đông Á}
\date{\today}

\begin{document}

\begin{frame}
    \titlepage
\end{frame}

\begin{frame}{Nội dung chương}
    \tableofcontents
\end{frame}

\section{1. Giới hạn của Deep Learning}

\begin{frame}{Bản chất thực sự của Deep Learning}
    \begin{columns}
        \column{0.5\textwidth}
        \begin{itemize}
            \item Mạng nơ-ron học sâu thực chất là những đường cong tham số khổng lồ (Parametric curves) được khớp (fit) với một lượng dữ liệu cực lớn.
            \item Mô hình không "suy nghĩ". Nó chỉ là một cơ sở dữ liệu \textbf{Nội suy tĩnh (Static Interpolative Database)}.
            \item Ở giai đoạn Inference, nó chỉ thực hiện truy xuất thông tin (Information retrieval) đã học.
        \end{itemize}
        
        \column{0.5\textwidth}
        \begin{figure}
            \centering
            \includegraphics[width=\textwidth,height=0.6\textheight,keepaspectratio]{../../images/ch19/ml_model.02e47549.png}
            \caption{Mô hình ML như một cơ sở dữ liệu nội suy}
        \end{figure}
    \end{columns}
\end{frame}

\begin{frame}{Sự mù lòa trước cái mới (Novelty)}
    \begin{columns}
        \column{0.5\textwidth}
        \begin{itemize}
            \item Điểm yếu cốt tử của LLMs: Rất giỏi làm những thứ từng xuất hiện trên mạng, nhưng bất lực trước một bài toán \textbf{Mới hoàn toàn}, dù rất đơn giản.
            \item Hình bên là một bài kiểm tra logic rất dễ với con người. Nhưng do không có trong dữ liệu huấn luyện, LLM không thể giải quyết được.
            \item "Năng lực" của LLM phụ thuộc vào \textit{Sự quen thuộc} chứ không phải \textit{Trí thông minh}.
        \end{itemize}
        
        \column{0.5\textwidth}
        \begin{figure}
            \centering
            \includegraphics[width=\textwidth,height=0.6\textheight,keepaspectratio]{../../images/ch19/arc_example_2.7d4f0c33.png}
            \caption{Một câu đố logic đơn giản nhưng hoàn toàn mới}
        \end{figure}
    \end{columns}
\end{frame}

\begin{frame}{Vấn đề "Học vẹt" ở Large Language Models}
    \begin{columns}
        \column{0.5\textwidth}
        \begin{itemize}
            \item Thử nghiệm với bài toán Monty Hall (Biến thể).
            \item Nếu bạn đổi một số chi tiết trong bài toán gốc (Ví dụ: Cửa trong suốt), LLM vẫn nhả ra câu trả lời kinh điển mà nó đã học vẹt.
            \item Điều này chứng tỏ mô hình không mô phỏng lại hoàn cảnh vật lý, nó chỉ \textbf{dự đoán token tiếp theo} khớp với dữ liệu phân bố của nó.
        \end{itemize}
        
        \column{0.5\textwidth}
        \begin{figure}
            \centering
            \includegraphics[width=\textwidth,height=0.6\textheight,keepaspectratio]{../../images/ch19/monty_hall.86dba4ca.png}
            \caption{Biến thể bài toán Monty Hall}
        \end{figure}
    \end{columns}
\end{frame}

\begin{frame}{Gót chân Achilles: Mẫu đối kháng (Adversarial Examples)}
    \begin{columns}
        \column{0.5\textwidth}
        \begin{itemize}
            \item Deep Learning rất dễ bị đánh lừa bởi \textbf{Adversarial Examples}.
            \item Thêm một lớp nhiễu rất nhỏ (Mắt người không nhận ra) vào hình ảnh con heo, AI sẽ tự tin 99\% đây là một chiếc máy bay.
            \item Nguyên nhân: AI không học cấu trúc vật lý của con heo, nó học các phân bố điểm ảnh cục bộ.
        \end{itemize}
        
        \column{0.5\textwidth}
        \begin{figure}
            \centering
            \includegraphics[width=\textwidth,height=0.6\textheight,keepaspectratio]{../../images/ch19/adversarial_example.8f3cfeb2.png}
            \caption{Đánh lừa AI bằng nhiễu ngẫu nhiên}
        \end{figure}
    \end{columns}
\end{frame}

\section{2. Bản chất Trí tuệ \& Thước đo ARC}

\begin{frame}{Trí tuệ là gì? (Skill vs Information)}
    \begin{columns}
        \column{0.5\textwidth}
        \begin{itemize}
            \item Việc nhớ toàn bộ Bách khoa toàn thư Wikipedia không làm bạn trở nên thông minh (Đó là \textbf{Information}).
            \item Trí tuệ thực sự là \textbf{Kỹ năng (Skill)} xử lý và thích nghi với các tình huống chưa từng gặp bao giờ.
            \item LLM hiện tại thiên về Information (Memorization) chứ chưa chạm đến Skill.
        \end{itemize}
        
        \column{0.5\textwidth}
        \begin{figure}
            \centering
            \includegraphics[width=\textwidth,height=0.6\textheight,keepaspectratio]{../../images/ch19/skill_vs_information.2468629a.png}
            \caption{Kỹ năng giải quyết vấn đề so với Khối lượng thông tin}
        \end{figure}
    \end{columns}
\end{frame}

\begin{frame}{Tổng quát hóa: Local vs Extreme Generalization}
    \begin{columns}
        \column{0.5\textwidth}
        \begin{itemize}
            \item Khả năng dự đoán các tình huống tương tự những gì đã học gọi là \textbf{Local Generalization} (Machine Learning hiện tại).
            \item Khả năng đối mặt với các tình huống không có tiền lệ (Ví dụ: hạ cánh khẩn cấp máy bay) gọi là \textbf{Extreme Generalization}.
            \item AGI (Artificial General Intelligence) đòi hỏi sự Tổng quát hóa cực hạn.
        \end{itemize}
        
        \column{0.5\textwidth}
        \begin{figure}
            \centering
            \includegraphics[width=\textwidth,height=0.6\textheight,keepaspectratio]{../../images/ch19/local_vs_extreme_generalization.4b71aa83.png}
            \caption{Mức độ Tổng quát hóa của AI so với Con người}
        \end{figure}
    \end{columns}
\end{frame}

\begin{frame}{Thước đo Trí tuệ: ARC Benchmark}
    \begin{columns}
        \column{0.5\textwidth}
        \begin{itemize}
            \item Francois Chollet đã tạo ra \textbf{Abstraction and Reasoning Corpus (ARC)} để đo lường khả năng Extreme Generalization của AI.
            \item ARC yêu cầu AI nhìn vào vài ví dụ (Input-Output) có tính logic trừu tượng, và suy ra lưới ảnh Output cuối cùng.
            \item Không thể giải bằng cách nhồi thêm dữ liệu (Scale up).
        \end{itemize}
        
        \column{0.5\textwidth}
        \begin{figure}
            \centering
            \includegraphics[width=\textwidth,height=0.6\textheight,keepaspectratio]{../../images/ch19/arc_example.2b17dc06.png}
            \caption{Một thử thách trong ARC Benchmark}
        \end{figure}
    \end{columns}
\end{frame}

\begin{frame}{AI SOTA vẫn thất bại trước ARC}
    \begin{columns}
        \column{0.5\textwidth}
        \begin{itemize}
            \item Ngay cả những mô hình mạnh nhất (OpenAI o3, Claude 3.5) cũng có tỷ lệ giải đúng ARC rất thấp so với con người.
            \item Điều này khẳng định triết lý: Việc mở rộng Tham số (Scaling Parameters) đang tiến tới mức bão hòa và không dẫn đến Trí tuệ thực thụ.
        \end{itemize}
        
        \column{0.5\textwidth}
        \begin{figure}
            \centering
            \includegraphics[width=\textwidth,height=0.6\textheight,keepaspectratio]{../../images/ch19/arc_task_not_solved_by_o3.03d72f08.png}
            \caption{Bài toán ARC mà các siêu AI vẫn đầu hàng}
        \end{figure}
    \end{columns}
\end{frame}

\begin{frame}{Sự phân cực trong ARC}
    \begin{columns}
        \column{0.5\textwidth}
        \begin{itemize}
            \item Có những bài toán con người mất 1 giây để giải, nhưng AI mất hàng triệu phép tính vẫn sai.
            \item Điều này cho thấy Kiến trúc tư duy của Deep Learning hoàn toàn khác biệt với Cognitive Architecture (Kiến trúc nhận thức) của con người.
        \end{itemize}
        
        \column{0.5\textwidth}
        \begin{figure}
            \centering
            \includegraphics[width=\textwidth,height=0.6\textheight,keepaspectratio]{../../images/ch19/arc_1_vs_arc_2.1201d48e.png}
            \caption{Sự sai biệt tư duy giữa AI và Não người}
        \end{figure}
    \end{columns}
\end{frame}

\section{3. Những Nền tảng còn thiếu của AI}

\begin{frame}{Bài học từ sinh học: Giun C. elegans}
    \begin{columns}
        \column{0.5\textwidth}
        \begin{itemize}
            \item Loài giun \textbf{C. elegans} chỉ có 302 tế bào thần kinh, nhưng nó có thể bò, tìm thức ăn, trốn kẻ thù và sinh sản.
            \item Trong khi mô hình AI có hàng trăm tỷ nơ-ron (Tham số) lại không thể tự điều khiển cơ thể linh hoạt trong thế giới thực.
            \item Sự sống là những hệ thống tự tổ chức (Self-organizing).
        \end{itemize}
        
        \column{0.5\textwidth}
        \begin{figure}
            \centering
            \includegraphics[width=\textwidth,height=0.6\textheight,keepaspectratio]{../../images/ch19/c_elegans.ca0de605.png}
            \caption{Sơ đồ Nơ-ron của giun C. elegans}
        \end{figure}
    \end{columns}
\end{frame}

\begin{frame}{Bùng nổ tổ hợp (Combinatorial Explosion)}
    \begin{columns}
        \column{0.5\textwidth}
        \begin{itemize}
            \item Thế giới thực không hữu hạn như bàn cờ vua.
            \item Giống như kính vạn hoa (Kaleidoscope), một vài mảnh kính nhỏ có thể tạo ra vô hạn các tổ hợp hình ảnh phức tạp.
            \item AI không thể đối phó với thế giới bằng cách "Học thuộc lòng" mọi tình huống.
        \end{itemize}
        
        \column{0.5\textwidth}
        \begin{figure}
            \centering
            \includegraphics[width=\textwidth,height=0.6\textheight,keepaspectratio]{../../images/ch19/kaleidoscope.fec15d2f.png}
            \caption{Sự đa dạng tổ hợp vô hạn}
        \end{figure}
    \end{columns}
\end{frame}

\begin{frame}{Value-Centric Abstraction (Khái niệm trung tâm)}
    \begin{columns}
        \column{0.5\textwidth}
        \begin{itemize}
            \item Deep Learning hiện tại hoạt động theo nguyên lý \textbf{Value-centric abstraction}.
            \item Trọng số (Weights) được tối ưu để tính ra một Giá trị (Value) làm đầu ra liên tục dựa trên các so sánh vector tĩnh.
            \item Nó không có khả năng thực thi các vòng lặp logic (While/For loops) vô định.
        \end{itemize}
        
        \column{0.5\textwidth}
        \begin{figure}
            \centering
            \includegraphics[width=\textwidth,height=0.6\textheight,keepaspectratio]{../../images/ch19/value_centric_abstraction.ea920558.png}
            \caption{Trừu tượng hóa theo Giá trị}
        \end{figure}
    \end{columns}
\end{frame}

\begin{frame}{Program-Centric Abstraction}
    \begin{columns}
        \column{0.5\textwidth}
        \begin{itemize}
            \item Con người suy luận bằng cách tạo ra các quy tắc logic (Algorithms).
            \item Tương lai của AI là \textbf{Program-centric abstraction}: Mô hình không xuất ra một Giá trị (Value), mà nó xuất ra một \textbf{Chương trình máy tính (Program)}.
        \end{itemize}
        
        \column{0.5\textwidth}
        \begin{figure}
            \centering
            \includegraphics[width=\textwidth,height=0.6\textheight,keepaspectratio]{../../images/ch19/program_centric_abstraction.028e5301.png}
            \caption{Trừu tượng hóa theo Chương trình}
        \end{figure}
    \end{columns}
\end{frame}

\section{4. Tương lai của AI: Meta-learning}

\begin{frame}{Tổng hợp Chương trình (Program Synthesis)}
    \begin{columns}
        \column{0.5\textwidth}
        \begin{itemize}
            \item Dò tìm trong một không gian các khối mã (Code blocks) để ghép lại thành một thuật toán.
            \item Đây là bước tiến tiềm năng để giải quyết ARC Benchmark, vì thuật toán là sự tổng quát hóa tối thượng (Extreme Generalization).
        \end{itemize}
        
        \column{0.5\textwidth}
        \begin{figure}
            \centering
            \includegraphics[width=\textwidth,height=0.6\textheight,keepaspectratio]{../../images/ch19/program_synthesis.8af8560a.png}
            \caption{AI tự tổng hợp chương trình để giải toán}
        \end{figure}
    \end{columns}
\end{frame}

\begin{frame}{Khái niệm Meta-learning (Học cách học)}
    \begin{columns}
        \column{0.5\textwidth}
        \begin{itemize}
            \item Trí tuệ con người không sinh ra đã biết làm toán. Con người có "Trí tuệ học tập".
            \item \textbf{Meta-learning} trong AI là việc huấn luyện một mô hình sao cho nó có khả năng \textit{tự cập nhật và tự học hỏi} thật nhanh trên môi trường mới chỉ với 1-2 ví dụ.
        \end{itemize}
        
        \column{0.5\textwidth}
        \begin{figure}
            \centering
            \includegraphics[width=\textwidth,height=0.6\textheight,keepaspectratio]{../../images/ch19/metalearning1.63aa6580.png}
            \caption{Nguyên lý của Meta-learning}
        \end{figure}
    \end{columns}
\end{frame}

\begin{frame}{Thuật toán Meta-learning (Inner \& Outer Loop)}
    \begin{columns}
        \column{0.5\textwidth}
        \begin{itemize}
            \item \textbf{Inner Loop:} Mô hình học cách giải quyết một tác vụ cụ thể nhanh nhất có thể.
            \item \textbf{Outer Loop:} Đạo hàm ngược qua cả quá trình học của Inner Loop để điều chỉnh siêu trọng số (Meta-weights), giúp hệ thống thích nghi nhanh hơn.
        \end{itemize}
        
        \column{0.5\textwidth}
        \begin{figure}
            \centering
            \includegraphics[width=\textwidth,height=0.6\textheight,keepaspectratio]{../../images/ch19/metalearning2.f5a5efde.png}
            \caption{Cơ chế Vòng lặp Học lồng nhau}
        \end{figure}
    \end{columns}
\end{frame}

\begin{frame}{AGI và Đích đến cuối cùng}
    \begin{itemize}
        \item \textbf{Artificial General Intelligence (AGI)} không phải là một LLM ghi nhớ mọi thứ trên Internet.
        \item AGI là một hệ thống có khả năng tự động tạo ra chương trình (Program Synthesis), áp dụng Meta-learning để thích nghi theo thời gian thực (Extreme Generalization), mà không phụ thuộc vào dữ liệu do con người mớm sẵn.
        \item Cuộc đua Trí tuệ Nhân tạo mới chỉ ở vạch xuất phát. Sinh viên hôm nay chính là những Kiến trúc sư của AGI ngày mai!
    \end{itemize}
\end{frame}

\end{document}
